\subsection{Observable Implications}

A new theoretical model should be able to account for many or even most of the established empirical findings that existing theories were able to explain. But any theory has its shortcomings and might perform better or worse in some respects than competing theories. Empirical research is then needed to find out which theories perform well and in which areas theories could be improved. This requires that theories be confronted with empirical evidence.

There are two meaningful ways to achieve this goal: testing implications where competing theories differ, or testing empirical implications of the model that it was not deliberately designed to reproduce \citep{cohenDevelopingSociologicalKnowledge1989, stinchcombeLogicSocialResearch2005}. This section aims to provide empirical researchers with initial implications of this model to assess it. The following hypotheses are predominantly derived from a separate computational experiment where the parameters are only varied around plausible values (analysis B in the online supplement). This list is a selection from the full set of hypotheses in the online supplement, and even that fuller list is not exhaustive. I encourage researchers to come up with their own hypotheses and, especially, to think about where this model differs from others.

One of the main motivations for the new theory is to link multiple phenomena in housing markets, most notably residential segregation, residential mobility, and housing inequality. As such, this model is able to produce new hypotheses about the relationships between these phenomena. As the main argument is that both households and landlords are spatially interdependent in their behavior, segregation is an important macro-context in which they act. Most hypotheses are based on this property of the model, making residential segregation a moderator for relationships between other variables. This choice also does not require empirically estimating the model parameters, which are often difficult to operationalize.

\begin{enumerate}
	\item lorem ipsum
	\item lorem ipsum
	\item lorem ipsum
\end{enumerate}